\documentclass[11pt]{article}
\usepackage[T1]{fontenc}
\usepackage[utf8]{inputenc}
\usepackage{amsmath,amssymb,amsthm}
\usepackage[margin=1in]{geometry}
\usepackage{array}
\usepackage[colorlinks=true,linkcolor=blue,citecolor=blue]{hyperref}
\usepackage{xcolor}

\newtheorem{theorem}{Theorem}[section]
\newtheorem{lemma}[theorem]{Lemma}
\newtheorem{proposition}[theorem]{Proposition}
\newtheorem{corollary}[theorem]{Corollary}
\newtheorem{assumption}[theorem]{Assumption}
\theoremstyle{definition}
\newtheorem{definition}[theorem]{Definition}
\theoremstyle{remark}
\newtheorem{remark}[theorem]{Remark}

\newcommand{\ASSUMED}{\textcolor{orange}{\textbf{[ASSUMED]}}}
\newcommand{\PROVED}{\textcolor{green!60!black}{\textbf{[PROVED]}}}
\newcommand{\IMPORTED}{\textcolor{blue}{\textbf{[IMPORTED]}}}
\newcommand{\Erec}{E_{\mathrm{rec}}}
\newcommand{\cA}{\mathcal{A}}
\newcommand{\cN}{\mathcal{N}}
\newcommand{\cM}{\mathcal{M}}
\newcommand{\cR}{\mathcal{R}}
\newcommand{\Sn}{\mathcal{S}_n}

\title{Split-regularized recoverability in Type III AQFT\\
\large Conditional expectations, split-dependent CMI, and an
audit-friendly recoverability contract}
\author{Lluis Eriksson}
\date{January 2026 (v1); July 2026 (v2)}

\begin{document}
\maketitle

\begin{abstract}
Local algebras in relativistic quantum field theory are typically Type
III, so reduced density matrices and von Neumann entropies are not
available without additional structure. We give a B-minimal,
audit-friendly interface for recoverability in Type III AQFT: we fix a
collar geometry and a split datum $N$ (an intermediate Type I factor)
and define (i) a split-regularized conditional mutual information (CMI)
and (ii) a Bures-fidelity-based recovery error for normal states. We
isolate, as explicit assumptions, the two hard bridges needed for an
exponential recoverability statement in Type III: (a) existence of an
$\omega_0$-preserving conditional expectation onto $N$ (a Takesaki-type
condition) and (b) an FR-type inequality in the fixed split
implementation. We prove a conditional theorem: if split-regularized CMI
decays exponentially in the collar width and an FR-type inequality holds
in that split, then recoverability error decays exponentially, with
constants tracked explicitly. This paper makes no Clay mass-gap claim
and does not invoke von Neumann entropy on Type III algebras without
split regularization. \emph{v2 (no v1 number is changed):} the
cross-reference labels are repaired (in v1 every assumption and theorem
was typeset as ``Definition~$x.y$'') and Table~\ref{tab:fals} is
re-typeset; a direction slip in the recovery candidate is corrected ---
under CE the GNS-adjoint $E_{\omega_0}^{\sharp\omega_0}$ is provably the
\emph{inclusion} $\iota: N \hookrightarrow \cM$, so the operational
candidate aligned with the recovery task is the predual
$(E_{\omega_0})_*$ (Remark~\ref{rem:direction}); the finite-dimensional
reduction is proved rather than remarked, including the equivalence of
the two split-regularized CMI definitions and $c_{\mathrm{FR}} = 1$ in
the fixed conventions (Lemma~\ref{lem:findim}); a constructive instance
of CE with product reference state is proved
(Proposition~\ref{prop:ce}); series positioning is added; and a
verification suite instantiates the entire contract end to end in the
Type I regime (gapped transverse-field Ising collar: CMI decay with
$\alpha \approx 1.06$, Petz-type reconstruction (the finite-dimensional Type I anchor) satisfying
$\Erec \le I^{(N)}$ at every width in the generated dataset, and the
$\omega_0$-adjoint identity $E^{\sharp\omega_0} = \iota$ at machine
precision).
\end{abstract}

\section{Scope and philosophy}

\textbf{Goal.} Provide a standalone, typed interface note for
recoverability in Type III AQFT that:
\begin{itemize}
\item avoids traces and finite-dimensional density matrices on Type III
algebras,
\item makes all dependence on the \emph{split datum} explicit,
\item separates proved content from bridge assumptions with
falsification routes.
\end{itemize}

\textbf{Non-goals.} We do not prove the Clay Yang--Mills mass gap and do
not claim that the split property alone implies exponential decay of a
Markovness proxy.

\textbf{Parameter convention.} Throughout, $\epsilon$ is a
\textbf{geometric} collar/separation parameter, while $\delta$ denotes
\textbf{tolerances} (regularization, thresholds). They are never
interchanged.

\section{Setup: regions, algebras, and normal states}

We work in a Haag--Kastler setting: to each spacetime region $O$ we
associate a von Neumann algebra $\cA(O)$ acting on a Hilbert space
$\mathcal{H}$, with isotony and locality. Fix three regions $A, B, C$
where $B$ is a collar separating $A$ from $C$ at scale $\epsilon$.

\begin{definition}[D1: tripartition at the algebra level]
\label{def:tripart}
Define the relevant local algebras
\[
\cA_A := \cA(A), \qquad \cA_B := \cA(B), \qquad \cA_{BC} := \cA(B \cup
C), \qquad \cA_{ABC} := \cA(A \cup B \cup C).
\]
A state $\omega$ is a \emph{normal} state on $\cA_{ABC}$, and $\omega_X$
denotes its restriction to $\cA_X$.
\end{definition}

\begin{remark}[Why we work with normal states]
\label{rem:normal}
Normal states are the natural notion of states on von Neumann algebras
in AQFT and interact well with modular theory, conditional expectations,
and Bures-type metrics.
\end{remark}

\section{Split datum (Type I implementation)}

\begin{assumption}[SPLIT: split inclusion datum]
\label{ass:split}
\ASSUMED. There exists a slightly thickened collar $B'$ with $B \subset
B' \subset B \cup C$ and an intermediate Type I factor $N$ such that
\[
\cA(B) \subset N \subset \cA(B') \subset \cA(B \cup C) = \cA_{BC}.
\]
We call $N$ the \emph{split datum} associated with the collar geometry
at scale $\epsilon$ \cite{DL}.
\end{assumption}

\begin{remark}[Typing intent]
\label{rem:typing}
SPLIT is not a decay statement. It is a structural input allowing a Type
I bookkeeping compatible with the collar geometry. All quantities
defined below are explicitly $N$-dependent.
\end{remark}

\begin{assumption}[SEL: split selection rule]
\label{ass:sel}
\ASSUMED. A selection rule $\epsilon \mapsto N(\epsilon)$ is fixed for
the regime of interest. All constants in decay statements are allowed to
depend on the chosen selection rule, but such dependence must be
declared.
\end{assumption}

\section{Channels and the recovery task (Heisenberg picture)}

\begin{remark}[Heisenberg picture convention]
\label{rem:heisenberg}
All channels between von Neumann algebras are taken in the Heisenberg
picture: a channel $\Phi: \cM_2 \to \cM_1$ is normal, completely
positive and unital. Its predual action on normal states is $\Phi_*:
\Sn(\cM_1) \to \Sn(\cM_2)$ given by $\Phi_*(\omega) := \omega \circ
\Phi$.
\end{remark}

\begin{definition}[D2: recovery channel and split-implemented recovery]
\label{def:recovery}
Fix a split datum $N$ for the collar geometry (as in
Assumption~\ref{ass:split}). A recovery channel is a normal CP unital
map (Heisenberg picture) $\cR: \cA_{BC} \to \cA_B$. Its predual map acts
on normal states by $\cR_*(\varphi) := \varphi \circ \cR$, hence $\cR_*:
\Sn(\cA_B) \to \Sn(\cA_{BC})$.

Given a normal state $\omega$ on $\cA_{ABC}$, the
\emph{split-implemented recovered state} (associated with the chosen
split datum $N$) is defined by applying $\cR_*$ on the $B$-leg \emph{in
the fixed split implementation induced by $N$}:
\[
\widetilde{\omega}^{(N)}_{ABC} := (\mathrm{id} \otimes \cR_*)\bigl(
\omega^{(N)}_{AB}\bigr),
\]
where $\omega^{(N)}_{AB}$ denotes the $AB$ marginal in that fixed split
implementation.
\end{definition}

\begin{remark}[Typing discipline]
\label{rem:discipline}
In Type III, $\cA(A \cup B)$ is not canonically a tensor product $\cA_A
\bar{\otimes} \cA_B$. Accordingly, the notation $(\mathrm{id} \otimes
\cR_*)(\omega^{(N)}_{AB})$ is shorthand for the standard tensor-leg
application performed \emph{inside the chosen split implementation
associated with $N$}. All dependence on this choice is tracked by the
superscript $(N)$.
\end{remark}

\section{Error metric in Type III: Bures fidelity and purified distance}

\begin{definition}[D3: fidelity/distance convention]
\label{def:fidelity}
Let $\omega, \varphi$ be normal states on a von Neumann algebra. We use
the \emph{squared} Bures fidelity $F_B(\omega, \varphi) \in [0,1]$ (so
$F_B = 1$ iff $\omega = \varphi$) \cite{Uhlmann}. As a convenient
monotone error parameter we adopt $P(\omega, \varphi) := \sqrt{1 -
F_B(\omega, \varphi)}$.
\end{definition}

\begin{definition}[D4: split-implemented recovery error functional]
\label{def:erec}
Fix a split datum $N$. Given $\omega$ on $\cA_{ABC}$ and a recovery
channel $\cR: \cA_{BC} \to \cA_B$, define
\[
\Erec(\omega; \cR; N) := -\log F_B\bigl(\omega_{ABC},
\widetilde{\omega}^{(N)}_{ABC}\bigr),
\]
where $\widetilde{\omega}^{(N)}_{ABC}$ is as in
Definition~\ref{def:recovery}.
\end{definition}

\begin{remark}[Convention discipline]
\label{rem:convention}
Different papers vary by using unsquared fidelity, Bures distance, or
$-2\log(\cdot)$. This note fixes Definitions~\ref{def:fidelity} and
\ref{def:erec}; any imported inequality must be aligned to these
conventions by an explicit constant (see Assumption~\ref{ass:frsplit};
in the fixed conventions the finite-dimensional FR inequality imports
with $c_{\mathrm{FR}} = 1$, Lemma~\ref{lem:findim}).
\end{remark}

\section{Conditional expectations and an Accardi--Cecchini recovery
candidate}

The Type III-specific ingredient is an $\omega_0$-preserving conditional
expectation onto the split factor.

\begin{assumption}[CE: $\omega_0$-preserving conditional expectation
onto $N$]
\label{ass:ce}
\ASSUMED. Fix a faithful normal reference state $\omega_0$ on $\cA_{BC}$
(e.g.\ a KMS state). Let $\cM := \cA_{BC}$, $\cN := N$ with $N$ the
split datum from Assumption~\ref{ass:split}. Assume there exists a
normal completely positive unital idempotent map $E_{\omega_0}: \cM \to
\cN$ such that $\omega_0 \circ E_{\omega_0} = \omega_0$.
\end{assumption}

\begin{remark}[Why CE is a genuine assumption]
\label{rem:cegenuine}
Existence of an $\omega_0$-preserving conditional expectation is tied to
modular invariance / Takesaki-type conditions and is not automatic
\cite{Takesaki}. This note isolates CE explicitly to prevent quantifier
drift. (A constructive Type I instance is proved in
Proposition~\ref{prop:ce}.)
\end{remark}

\begin{definition}[D5: GNS-adjoint w.r.t.\ $\omega_0$]
\label{def:gnsadjoint}
Let $\Phi: \cM \to \cN$ be normal CP and let $\omega_0$ be faithful.
Define $\langle X, Y \rangle_{\omega_0} := \omega_0(X^* Y)$. When it
exists, the $\omega_0$-adjoint $\Phi^{\sharp\omega_0}: \cN \to \cM$ is
defined by $\langle \Phi(X), Y \rangle_{\omega_0} = \langle X,
\Phi^{\sharp\omega_0}(Y) \rangle_{\omega_0}$ for all $X \in \cM$, $Y \in
\cN$.
\end{definition}

\begin{definition}[D6: Accardi--Cecchini / Petz-dual candidate]
\label{def:ac}
Assume CE and that the $\omega_0$-adjoint exists for $E_{\omega_0}$.
Define the canonical recovery candidate $\cR^{\mathrm{AC}} :=
E_{\omega_0}^{\sharp\omega_0}: \cN \to \cM$.
\end{definition}

\begin{remark}[Candidate vs existential recovery]
\label{rem:candidate}
This note does not claim that $\cR^{\mathrm{AC}}$ achieves an FR-type
bound pointwise; it provides a canonical, typed candidate under CE.
Proving optimality or near-optimality of $\cR^{\mathrm{AC}}$ is a
separate problem. (See Remark~\ref{rem:direction} for the v2 correction
of the candidate's direction.)
\end{remark}

\begin{remark}[v2 correction: $E^{\sharp\omega_0} = \iota$, and the
operational candidate is the predual]
\label{rem:direction}
\PROVED\ (v2). Under CE the GNS-adjoint of Definition~\ref{def:ac} can
be computed exactly: for $X \in \cM$, $Y \in \cN$, the module property
of conditional expectations and $\omega_0$-invariance give
\[
\langle E_{\omega_0}(X), Y\rangle_{\omega_0}
= \omega_0\bigl(E_{\omega_0}(X)^* Y\bigr)
= \omega_0\bigl(E_{\omega_0}(X^* \iota(Y))\bigr)
= \omega_0\bigl(X^* \iota(Y)\bigr)
= \langle X, \iota(Y)\rangle_{\omega_0},
\]
so $E_{\omega_0}^{\sharp\omega_0} = \iota$, the \emph{inclusion} $\cN
\hookrightarrow \cM$. Consequently the predual of
$\cR^{\mathrm{AC}} = \iota$ is the \emph{restriction} $\Sn(\cM) \to
\Sn(\cN)$ --- the wrong direction for the recovery task of
Definition~\ref{def:recovery}. The operational recovery candidate
aligned with D2 is the \emph{predual of the conditional expectation
itself}, $(E_{\omega_0})_*: \Sn(\cN) \to \Sn(\cM)$ (reattachment), while
the pair $(E_{\omega_0}, \iota)$ is precisely the GNS-adjoint duality
pair. This matches the corrected trace-predual duality of the companion
2601.0046 (v2), where the same direction discipline was established for
the Accardi--Cecchini generalized conditional expectation. The
verification suite checks $E^{\sharp\omega_0} = \iota$ and the explicit
reattachment form of $(E_{\omega_0})_*$ at machine precision in the Type
I instance of Proposition~\ref{prop:ce}.
\end{remark}

\section{Split-regularized conditional mutual information}

We define a split-regularized CMI $I^{(N)}_\omega(A : C | B)$ that is
meaningful given a split datum $N$.

\subsection{Two definitions (kept explicitly $N$-dependent)}

\begin{definition}[D7a: split-implemented CMI (Type I bookkeeping)]
\label{def:cmia}
Assume SPLIT and fix $N$. Use the induced Type I implementation to
represent the relevant inclusions in a tensor-product bookkeeping
compatible with $(A, B, C)$. Define $I^{(N)}_\omega(A : C | B)$ by
applying the usual Type I CMI formula inside that implementation (with
the conventions fixed by the chosen implementation), allowing the value
$+\infty$ if the expression diverges.
\end{definition}

\begin{definition}[D7b: Araki-relative-entropy CMI (alternative)]
\label{def:cmib}
Alternatively, define $I^{(N)}_\omega(A : C | B)$ using Araki relative
entropy in standard form on the relevant inclusions, with explicit
dependence on the same split datum $N$ and the chosen inclusions.
\end{definition}

\begin{remark}[No equivalence claimed in general]
\label{rem:noequiv}
We do not assume Definition~\ref{def:cmia} and Definition~\ref{def:cmib}
coincide in full generality. Any equivalence statement must be proved or
explicitly assumed (see Appendix~\ref{app:equiv}; the finite-dimensional
case is \emph{proved} in v2, Lemma~\ref{lem:findim}).
\end{remark}

\section{Main contract: conditional exponential recoverability}

\begin{assumption}[CMI-DECAY: split-regularized exponential Markov
bound]
\label{ass:cmidecay}
\ASSUMED. Assume Assumption~\ref{ass:sel} and work with the selected
split data $N(\epsilon)$. Fix a family of normal states
$\{\omega^{(\epsilon)}\}_\epsilon$ on $\cA_{ABC}$ (or a declared state
class, stated explicitly). Assume there exist constants $K > 0$ and
$\alpha > 0$ such that for all $\epsilon$ in the regime of interest,
\[
I^{(N(\epsilon))}_{\omega^{(\epsilon)}}(A : C | B) \le K e^{-\alpha
\epsilon}.
\]
Constants may depend on the model, the state family/class, and the
selection rule $\epsilon \mapsto N(\epsilon)$, but not on $\epsilon$
itself.
\end{assumption}

\begin{assumption}[FR-SPLIT: FR-type inequality in the fixed split]
\label{ass:frsplit}
\ASSUMED\ (import/bridge). Fix the split datum $N$ and a choice of
split-regularized CMI definition (either Definition~\ref{def:cmia} or
Definition~\ref{def:cmib}). Assume there exists a universal constant
$c_{\mathrm{FR}} > 0$ (depending only on the conventions
Definitions~\ref{def:fidelity} and \ref{def:erec} and the chosen CMI
definition) such that for every normal state $\omega$ on $\cA_{ABC}$
there exists a recovery channel $\cR: \cA_{BC} \to \cA_B$ (in the fixed
split implementation) satisfying
\[
\Erec(\omega; \cR; N) \le c_{\mathrm{FR}}\, I^{(N)}_\omega(A : C | B).
\]
\end{assumption}

\begin{theorem}[T1: conditional exponential recoverability (Type III,
split-regularized)]
\label{thm:contract}
\PROVED\ (given Assumptions~\ref{ass:cmidecay} and \ref{ass:frsplit}).
Under Assumptions~\ref{ass:cmidecay} and \ref{ass:frsplit}, for each
$\epsilon$ there exists a recovery channel $\cR$ such that
\[
\Erec(\omega; \cR; N(\epsilon)) \le c_{\mathrm{FR}}\, K\,
e^{-\alpha\epsilon}.
\]
Equivalently, with $P$ as in Definition~\ref{def:fidelity},
$P\bigl(\omega_{ABC}, \widetilde{\omega}^{(N(\epsilon))}_{ABC}\bigr) \le
\sqrt{1 - \exp(-c_{\mathrm{FR}} K e^{-\alpha\epsilon})}$.
\end{theorem}

\begin{proof}
Combine Assumption~\ref{ass:cmidecay} with Assumption~\ref{ass:frsplit}
and the definition of $\Erec$ in Definition~\ref{def:erec}. The
purified-distance form follows from Definition~\ref{def:fidelity}.
\end{proof}

\begin{remark}[Where the real work is]
\label{rem:realwork}
T1 is a contract statement. The substantive work in applications is: (i)
verifying CMI-DECAY for a concrete AQFT model and a concrete selection
rule $\epsilon \mapsto N(\epsilon)$, and (ii) importing/proving an
FR-type inequality in the fixed split implementation aligned to the
conventions Definitions~\ref{def:fidelity} and \ref{def:erec}.
\end{remark}

\section{Falsification matrix (minimal)}

\begin{table}[h!]
\centering
\begin{tabular}{p{3.0cm} p{2.2cm} p{4.2cm} p{4.6cm}}
\hline
\textbf{Item} & \textbf{Status} & \textbf{Verification proxy} &
\textbf{Concrete falsifier} \\
\hline
SPLIT ($N(\epsilon)$ exists) & \ASSUMED & cite conditions; construct
$N(\epsilon)$ & exhibit model/geometry where split fails \\
SEL (selection rule fixed) & \ASSUMED & specify $N(\epsilon)$ explicitly
& show ambiguity changes constants qualitatively \\
CE ($E_{\omega_0}$ exists) & \ASSUMED\ general; instance \PROVED\ (v2) &
check modular invariance / Takesaki condition & exhibit Takesaki
obstruction / lack of $\omega_0$-invariance \\
CMI-DECAY & \ASSUMED & estimate $I^{(N)}_\omega(A:C|B)$ vs $\epsilon$ &
show non-decay or slower-than-claimed decay \\
FR-SPLIT & \ASSUMED\ general; Type I import \PROVED\ (v2) & import a
theorem aligned to Definitions~\ref{def:fidelity} and \ref{def:erec} &
give a counterexample under stated definitions \\
\hline
\end{tabular}
\caption{Minimal falsification matrix for the Type III recoverability
interface (re-typeset in v2; v1's table had overlapping column text).}
\label{tab:fals}
\end{table}

\section{Compatibility with finite-dimensional benchmarks}

\begin{remark}[Finite-dimensional reduction]
\label{rem:findim}
When the relevant algebras are Type I (finite-dimensional), the split
datum can be taken canonically and split-regularized CMI reduces to the
usual finite-dimensional CMI. In that regime, FR-type inequalities
reduce to their standard QIT forms, up to the constant $c_{\mathrm{FR}}$
determined by conventions. (Proved in v2 as Lemma~\ref{lem:findim}.)
\end{remark}

\begin{lemma}[Finite-dimensional reduction, proved form (v2)]
\label{lem:findim}
\PROVED. Let $\cA_{ABC} = B(\mathcal{H}_A \otimes \mathcal{H}_B \otimes
\mathcal{H}_C)$ with all factors finite-dimensional, and take the
canonical split datum $N = B(\mathcal{H}_B)$. Then: (i)
Definition~\ref{def:cmia} and Definition~\ref{def:cmib} coincide:
\[
S_{AB} + S_{BC} - S_B - S_{ABC}
= D\bigl(\rho_{ABC} \,\|\, \rho_A \otimes \rho_{BC}\bigr)
- D\bigl(\rho_{AB} \,\|\, \rho_A \otimes \rho_B\bigr),
\]
where in finite dimension Araki relative entropy equals the standard
$D(\rho\|\sigma) = \mathrm{Tr}\,\rho(\log\rho - \log\sigma)$. (ii) In
the conventions of Definitions~\ref{def:fidelity} and \ref{def:erec}
(squared Bures fidelity, $\Erec = -\log F_B$), the Fawzi--Renner theorem
\cite{FR} imports as FR-SPLIT with $c_{\mathrm{FR}} = 1$: there exists a
recovery channel (a rotated Petz map) with
$\Erec \le I(A:C|B)$.
\end{lemma}

\begin{proof}
(i) Expanding both relative entropies with marginals of the same global
state, $D(\rho_{ABC}\|\rho_A \otimes \rho_{BC}) = S_A + S_{BC} -
S_{ABC}$ (mutual information $I(A:BC)$) and $D(\rho_{AB}\|\rho_A \otimes
\rho_B) = S_A + S_B - S_{AB}$ ($I(A:B)$); the difference is the chain
rule $I(A:BC) - I(A:B) = I(A:C|B)$, which equals the entropy
combination. In finite dimension Araki's definition reduces to the
standard trace formula, so D7b built on these inclusions is the
right-hand side. (ii) Fawzi--Renner \cite{FR} states $I(A:C|B) \ge
-\log F$ with $F$ the squared fidelity of recovery for a rotated Petz
channel; since Definition~\ref{def:erec} uses $\Erec = -\log F_B$ with
squared $F_B$, the import constant is exactly $1$. (This is the
convention discipline of Remark~\ref{rem:convention}; with unsquared
fidelity the constant would be $2$, the recurring factor tracked across
this series.)
\end{proof}

\begin{proposition}[A constructive CE instance (v2)]
\label{prop:ce}
\PROVED. Let $\cM = B(\mathcal{H}_N \otimes \mathcal{H}_K)$, $\cN =
B(\mathcal{H}_N) \otimes \mathbf{1}_K$, and let $\omega_0 =
\mathrm{Tr}[(\rho_N \otimes \rho_K)\,\cdot\,]$ be a faithful product
state. Then $E_{\omega_0}(X) := \mathrm{Tr}_K[(\mathbf{1} \otimes
\rho_K) X] \otimes \mathbf{1}_K$ is a normal CP unital idempotent map
onto $\cN$ with $\omega_0 \circ E_{\omega_0} = \omega_0$: the Takesaki
condition is satisfied constructively. Its GNS-adjoint is the inclusion
(Remark~\ref{rem:direction}), and its predual acts as the explicit
reattachment $(E_{\omega_0})_*(\varphi_N) = \varphi_N \otimes \rho_K$.
\end{proposition}

\begin{proof}
Unitality, idempotence and complete positivity are immediate from the
form of $E_{\omega_0}$ (partial trace against a state followed by tensor
extension). $\omega_0$-invariance:
$\omega_0(E_{\omega_0}(X)) = \mathrm{Tr}[\rho_N\,
\mathrm{Tr}_K((\mathbf{1} \otimes \rho_K)X)] = \mathrm{Tr}[(\rho_N
\otimes \rho_K) X] = \omega_0(X)$. The predual form follows by pairing
$\varphi_N \circ E_{\omega_0}$ against arbitrary $X$. The suite verifies
all five properties numerically (Choi positivity to $10^{-16}$).
\end{proof}

\begin{remark}[Series positioning (v2)]
\label{rem:series}
This note is the AQFT companion referenced as ``forthcoming'' across the
series. The two CMI definitions correspond to the series'
split-regularized variants: D7a is the embedded (Type I bookkeeping)
CMI of 2601.0020, D7b is the Araki-difference CMI of 2601.0034, and the
split-implementation discipline follows 2601.0046, whose corrected
trace-predual duality is what Remark~\ref{rem:direction} instantiates
for CE. Finite-dimensional FR/Petz numerics aligned with these
conventions appear in 2601.0007 and 2601.0035; collar-geometry CMI
decay and its saturation structure in 2601.0050. The verification suite
(\texttt{verification/2601-0065/}) instantiates the full contract in the
Type I regime: a gapped transverse-field Ising collar ($g = 1.5$, $A$
fixed at $2$ sites, $C$ at $4$, collar width $w = 1..4$) gives
$I^{(N)} \approx 0.20\, e^{-1.06 w}$, and the Petz-type/CE-pullback reconstruction channel
satisfies $\Erec \le I^{(N)}$ at every width in the generated dataset,
i.e.\ the chain of Theorem~\ref{thm:contract} holds end to end with
$c_{\mathrm{FR}} = 1$.
\end{remark}

\section*{Note added in v2}
No v1 definition, assumption, theorem statement or numbering is changed;
new v2 material is appended at the end of the affected sections. v2: (1)
repairs the cross-reference labels --- in v1 every reference was typeset
as ``Definition~$x.y$'' regardless of environment, so e.g.\ ``given
Definitions 8.1 and 8.2'' in v1 denotes
Assumptions~\ref{ass:cmidecay} and \ref{ass:frsplit}, and ``as in
Definition 3.1'' denotes Assumption~\ref{ass:split}; (2) re-typesets
Table~\ref{tab:fals} (overlapping, partially illegible entries in v1);
(3) corrects a direction slip in the recovery candidate: under CE the
GNS-adjoint $E_{\omega_0}^{\sharp\omega_0}$ is provably the inclusion
$\iota$, so the predual of $\cR^{\mathrm{AC}}$ as defined in v1 is a
restriction, not a recovery; the operational candidate is
$(E_{\omega_0})_*$ (Remark~\ref{rem:direction}, verified at machine
precision); (4) proves the finite-dimensional reduction, including
D7a $=$ D7b and the $c_{\mathrm{FR}} = 1$ import of Fawzi--Renner in the
fixed conventions (Lemma~\ref{lem:findim}), settling
Appendix~\ref{app:equiv} in the Type I regime; (5) proves a constructive
CE instance with product reference state (Proposition~\ref{prop:ce});
(6) adds series positioning (v1 cited no companion papers;
Remark~\ref{rem:series}) and the verification suite instantiating the
entire contract end to end in the Type I regime.

\appendix

\section{Optional: equivalence of CMI definitions}
\label{app:equiv}

If one wants to identify Definition~\ref{def:cmia} and
Definition~\ref{def:cmib} (for a fixed choice of inclusions and split
data), this requires an explicit theorem-level statement. In full
(Type III) generality this paper does not assume such equivalence; it
should be treated as \IMPORTED\ (with hypotheses restated verbatim) or
as \ASSUMED\ (with a falsification route). \emph{v2:} in the
finite-dimensional (Type I) regime the equivalence is \PROVED\
(Lemma~\ref{lem:findim}(i)) and verified numerically to machine
precision.

\section{Reference-state dependence and model dependence}
\label{app:refstate}

The CE assumption (Assumption~\ref{ass:ce}) depends on the choice of
faithful normal reference state $\omega_0$. In AQFT applications, a
natural choice is often a KMS state. Any claim using CE should
explicitly record which $\omega_0$ is used and whether $\omega_0$ varies
with $\epsilon$.

\begin{thebibliography}{12}
\bibitem{DL} S. Doplicher and R. Longo, \emph{Standard and split
inclusions of von Neumann algebras}, Invent. Math. \textbf{75}, 493--536
(1984).
\bibitem{Takesaki} M. Takesaki, \emph{Theory of Operator Algebras I},
Springer, 1979.
\bibitem{Petz} D. Petz, \emph{Sufficient subalgebras and the relative
entropy of states of a von Neumann algebra}, Commun. Math. Phys.
\textbf{105}, 123--131 (1986).
\bibitem{FR} O. Fawzi and R. Renner, \emph{Quantum conditional mutual
information and approximate Markov chains}, Commun. Math. Phys.
\textbf{340}, 575--611 (2015).
\bibitem{Uhlmann} A. Uhlmann, \emph{The ``transition probability'' in
the state space of a $*$-algebra}, Rep. Math. Phys. \textbf{9}, 273--279
(1976).
\bibitem{STH} D. Sutter, M. Tomamichel, and P. Harremo\"es,
\emph{Strengthened monotonicity of relative entropy via pinched Petz
recovery map}, IEEE Trans. Inf. Theory \textbf{62}, 2907--2913 (2016).
\bibitem{S0046} L. Eriksson, ai.viXra:2601.0046 (2026; v2 2026).
\bibitem{S0034} L. Eriksson, ai.viXra:2601.0034 (2026; v2 2026).
\bibitem{S0020} L. Eriksson, ai.viXra:2601.0020 (2026; v2 2026).
\bibitem{S0035} L. Eriksson, ai.viXra:2601.0035 (2026; v2 2026).
\bibitem{S0007} L. Eriksson, ai.viXra:2601.0007 (2026; v2 2026).
\bibitem{S0050} L. Eriksson, ai.viXra:2601.0050 (2026; v2 2026).
\end{thebibliography}

\end{document}
